% profiles/university-of-california-berkeley.tex
% — University of California-Berkeley
% ============================================================================
% source: https://grad.berkeley.edu/academics/degree-progress/dissertation/
%         (Berkeley Graduate Division, "Dissertation Writing and Filing" — the
%          filing guide itself, carrying the Accessibility, Formatting your
%          manuscript, Special Page Formats and Organizing your manuscript
%          sections quoted throughout this file)
% source: https://docs.google.com/document/d/1qGrxd-rnFRasK1HTDLfQOT0Y_X6iSFx9ni1xJRIDrUg/
%         (the Graduate Division's own "Sample Title Page", linked from that
%          page. Read both as text and as a rendered page. It carries its own
%          authoritative "Notes on How to Have an Effective Title Page", quoted
%          below.)
% accessed: 2026-09-05
% checked:  texlib-thesis-profile-round
%
% VERIFIED and set here: geometry, spacing, the title page, the absence of an
% approval page, and two accessibility check-flags. NOT set: chapter-opening
% margin (no such rule).
%
% ============================================================================
% BERKELEY REQUIRES SINGLE SPACING. READ THAT AGAIN BEFORE EDITING THIS FILE.
%
%   "Spacing: Your manuscript must be single-spaced throughout, including the
%    abstract, dedication, acknowledgments, and introduction."
%
% Every other institution in this directory either requires double spacing or
% permits it. Berkeley requires the opposite, for the whole manuscript, and says
% so flatly. A well-meaning pass that "fixes" \thesissetspacing{single} to
% {double} for consistency with its neighbours would put a filer out of spec on
% every page. This is exactly the cross-contamination the round instructions
% warn about, and it is why the warning is at the top of the file rather than
% beside the setting.
% ============================================================================

\thesisinstitution{University of California, Berkeley}

% --- Filing figures ----------------------------------------------------------
% "Margins: For the manuscript material, including headers, footers, tables,
% illustrations, and photographs, all margins must be at least 1 inch from the
% edges of the paper. Page numbers must be 3/4 of an inch from the edge."
%
% A MINIMUM, so 1in is the floor and is set. Note how wide the scope is —
% headers, footers, tables, illustrations and photographs are all named, and
% elsewhere: "Tables, charts, and graphs ... must fit within the required
% margins", with permission to "reduce the size of a page to fit within the
% required margins".
%
% The page-number figure is also a minimum, stated in the pagination section as
% "must be positioned either in the upper right corner, lower right corner, or
% the bottom center and must be at least 3/4 of an inch from the edges". This
% class's `includefoot' geometry keeps the number inside the 1in text block,
% i.e. at least 1in from the edge, which satisfies "at least 3/4 of an inch"
% with room to spare. Conservative, and in spec.
\thesissetgeometry{left=1in, right=1in, top=1in, bottom=1in}

% See the banner above. This is a REQUIREMENT, and it is single.
\thesissetspacing{single}

% No chapter-opening rule is published, so \thesissetchapteropening stays unset.

% --- Accessibility -----------------------------------------------------------
% THE MOST EXPLICIT ACCESSIBILITY MANDATE OF ANY INSTITUTION IN THIS DIRECTORY,
% and unlike Florida's it is stated entirely in the graduate school's own prose
% rather than read off a template:
%
%   "As of April 2026, all electronic theses and dissertations (ETDs) must
%    conform to WCAG 2.1 Level AA web accessibility requirements before
%    submission to comply with nondiscrimination provisions outlined in Title II
%    of the Americans with Disabilities Act."
%
% NOTE THE DATE: April 2026, the large-public-entity ADA Title II deadline.
% Other profiles in this directory legitimately carry April 2027, which is the
% date for smaller public entities. Both are right. Do not reconcile them.
%
% \thesisrequiretagging is raised on: "PDFs should contain code or markup
% language (also known as tags) that ensures that the content is navigable and
% readable by screen readers or other assistive technology."
%
% \thesisrequiredocumentlanguage is raised on the Key Components list:
% "Metadata (file properties): Add document description such as author, title,
% and language."
%
% \thesisrequirepdfstandard is NOT raised, deliberately. Berkeley names WCAG 2.1
% Level AA, which is a WEB CONTENT standard, not a PDF standard — it is not the
% kind of thing \GetDocumentProperties{document/pdfstandard} reports, and
% declaring it there would make the class compare two incomparable strings and
% emit an advisory that could not be acted on. It is recorded verbatim below
% instead. A reviewer who wants the build-time nag anyway should say so.
\thesisrequiretagging
\thesisrequiredocumentlanguage
\thesisaccessibilityrequirement{WCAG 2.1 Level AA}
	{"As of April 2026, all electronic theses and dissertations (ETDs) must
	 conform to WCAG 2.1 Level AA web accessibility requirements before
	 submission to comply with nondiscrimination provisions outlined in Title II
	 of the Americans with Disabilities Act."}
\thesisaccessibilityrequirement{Tags}
	{"PDFs should contain code or markup language (also known as tags) that
	 ensures that the content is navigable and readable by screen readers or
	 other assistive technology." Tags named: "Headings, Lists, Links, Tables,
	 Paragraphs, and other document elements".}
\thesisaccessibilityrequirement{Tables}
	{"PDF must have table headers for columns and rows."}
\thesisaccessibilityrequirement{Reading order}
	{"Ensure that tags and headings are ordered to support the logical reading
	 order."}
\thesisaccessibilityrequirement{Alt text}
	{"Ensures that all images, figures, or other visual elements are described."
	 Berkeley points at the Section 508 guide to authoring meaningful
	 alternative text.}
\thesisaccessibilityrequirement{Colour}
	{"Choose colors with sufficient contrast, and do not rely on color alone to
	 convey information."}
\thesisaccessibilityrequirement{Metadata}
	{"Add document description such as author, title, and language."}
\thesisaccessibilityrequirement{Whose job it is}
	{"It is the responsibility of the students to ensure their ETDs meet
	 accessibility requirements. Graduate Division staff are not experts in
	 manuscript accessibility formatting and authoring tools." Berkeley expects
	 the manuscript to pass the Adobe Acrobat Pro Accessibility Checker, and
	 says its reviewers use that tool.}

% --- Title page --------------------------------------------------------------
% Reproduced from the Graduate Division's own Sample Title Page. Its notes are
% prescriptive and are followed here:
%
%   * "Line Breaks: Line breaks must match this example." The eight-line block
%     below is the sample's, verbatim, and MUST NOT be reflowed — note that
%     "in", "in the" and "of the" each occupy a line of their own, which looks
%     like an accident and is not.
%   * "Title: Use full title of dissertation, must be in mixed capitalization."
%     So the title is emitted verbatim — not uppercased.
%   * "Do not bold any text on your title page" (filing guide, Special Page
%     Formats). Nothing here is bold, unlike the class's neutral page.
%   * "The title page does not contain page numbers."
%   * "Name: You must use your registered student name as it is currently
%     recorded with the Office of the Registrar."
%   * "Degree: List the specific degree you are filing to receive, spelled out
%     in full. Do not use abbreviations (e.g., list Doctor of Philosophy, not
%     Ph.D.). Do not list previous degrees on your title page." — so \author
%     carries the name ALONE here, with no ", B.S." tail. That is the opposite
%     of Emory and MD Anderson in this same round.
%   * "University: Write out fully 'University of California, Berkeley' do not
%     abbreviate to UC Berkeley or just University of California." Hard-coded
%     below for that reason.
%   * "Major: List the major in which your degree will be awarded. Do not list
%     your specialization or area of emphasis." Also: "All majors in the College
%     of Engineering, except Bioengineering and Computer Science, must put
%     Engineering- or Engineering Science- before the name of the major."
%     A document supplies that through \gradprogram.
%   * "Committee members: List the committee in charge of your manuscript, with
%     the chair or Co-Chairs listed first. If you have Co-Chairs, use a separate
%     line for each name. Co-Chair must be hyphenated and capitalized exactly as
%     'Co-Chair'." The sample prefixes each name with "Professor", so a document
%     should write \committeemember{Professor First I. Last}{Co-Chair}.
%   * "Semester and year: Degrees are conferred in Fall, Spring, and Summer.
%     List the semester and year in which your degree will be conferred."
%     So \graddate{Spring 2027} — a SEMESTER, not a month. Berkeley is the only
%     institution in this round that wants a semester name there.
\renewcommand{\thesistitlepage}{%
	\loadgeometry{nopagenumber}%
	\begin{titlepage}%
		\thispagestyle{empty}\singlespacing\centering
		\vspace*{0.5in}
		\thesis@tagtitle{Title}{\@title}\par
		\vspace{2\baselineskip}
		By\par
		\vspace{\baselineskip}
		\@author\par
		\vspace{3\baselineskip}
		% The sample's line breaks, verbatim. Do not reflow.
		A \thesis@docnoun{} submitted in partial satisfaction of the\\
		requirements for the degree of\\
		\thesis@degree\\
		in\\
		\thesis@program\\
		in the\\
		Graduate Division\\
		of the\\
		University of California, Berkeley\par
		\vspace{3\baselineskip}
		Committee in charge:\par
		\vspace{\baselineskip}
		\begingroup
		% "Professor First I. Last, Co-Chair" -- role after the name, one per
		% line, no rules.
		\renewcommand{\thesis@memberline}[2]{%
			##1\ifx\relax##2\relax\else, ##2\fi\par
		}%
		\thesis@committee
		\endgroup
		\vspace{3\baselineskip}
		\thesis@date\par
		\vspace*{\fill}
	\end{titlepage}%
	\loadgeometry{normal}%
}

% --- Committee approval page --------------------------------------------------
% Berkeley's filed manuscript has none. "Organizing your manuscript" gives the
% page order in full — "Title Page, Copyright page or a blank page, Abstract,
% Optional preliminary pages such as Dedication page, Table of contents, List of
% figures ..., Preface or introduction, Acknowledgments, Curriculum Vitae, Main
% text, References, Bibliography, Appendices" — and contains no approval,
% committee or signature page. The committee is named on the title page above.
%
% PHYSICAL SIGNATURES: none, anywhere. Approval is collected through the FINAL
% SIGNATURE eFORM in CalCentral, which is step 1 of filing and lives entirely
% outside the manuscript. Berkeley cross-checks it against the title page: "the
% committee listed on your title page (and on the final signature eForm you will
% submit) must match your currently approved committee reflected in your
% CalCentral" record — so the committee list above is not decoration, and a
% mismatch will stop a filing.
\thesisnoapprovalpage
